\section{Introduction}

A prime-row collision is useful only when it belongs to the literal source support.
This distinction becomes decisive immediately after the cutoff-one regime.  The
preceding audit proved that the clock $H=4Q^2$, $h=4Q$ has pairwise disjoint prime
supports, so its AP energy equals its diagonal energy exactly \cite{tpc225}.  A natural
response is to enlarge the modulus and admit more multiplier layers.  The first apparent
overlap, however, uses a nonprimitive multiplier and disappears once the TPC row is
restored.

The broader analytic motivation comes from the role of primes in arithmetic progressions
and large-sieve estimates in sieve bounds related to twin primes
\cite{lichtman2023primes}.  The present paper does not import such an estimate.  It asks
a prior structural question: when does the exact finite row geometry first permit a
cross-prime term, and what determines its sign?  The answer separates collision
existence from cancellation.

We study the integer-dilated family
\[
  H=4Q^2,\qquad h_L=4LQ,\qquad L=1,2,3,4,
\]
while retaining the primitive multiplier condition, the common normalization
$C_h=1/h_L$, and the four-packet energy interface of TPC-220 and TPC-224
\cite{tpc220,tpc224}.  The paper makes three falsifiable contributions.

\begin{enumerate}[leftmargin=2em]
  \item We prove pairwise disjointness for every $L\le3$ in the stable range
  $Q\ge8$.  This rejects the tempting $L=3$, $m=4$ overlap as a nonprimitive
  artifact.
  \item We classify the first primitive transition.  At $L=4$, all collisions have
  the form $7p+3r=16Q$ and use only the multiplier pair $(3,-7)$, together with its
  global negative.
  \item We prove that the first-transition Gram correction is sign-indefinite across
  profile families.  Constant and inherited affine profiles amplify the AP energy,
  whereas balanced sign profiles produce strict AP saving and complete packet
  cancellation.
\end{enumerate}

The strongest finite witness occurs at $Q=25$.  The same full prime shell gives
$\EAP/\Ediag=15/13$ for aligned profiles and $11/13$ for balanced sign profiles.
Thus overlap is necessary for the correction, but overlap alone does not pay an AP
saving.  The missing theorem must derive a negative sign from the physical packet
source rather than select it as a finite profile.
